%! Singular Values and Singular Vectors — Unit 7, Lesson 7.1 — Exit Ticket
\documentclass[10pt]{article}
\usepackage{linalg-article}
\usepackage{linalg-boxes}
\newcommand{\TallMath}[1]{$\displaystyle #1\rule[-1.4em]{0pt}{3.2em}$}
\newcommand{\uu}{\mathbf{u}}
\newcommand{\vv}{\mathbf{v}}
\newcommand{\T}{^{\top}}

\begin{document}

\pageheader{Unit 7, Lesson 7.1}{Exit Ticket}
\namedateperiod

\vspace{0.15in}

No notes. Show your reasoning. Use $A=\begin{bmatrix} 1 & -2 \\ 2 & 2 \end{bmatrix}$ for Questions 1--2.

\begin{enumerate}[leftmargin=1.8em, itemsep=16pt]
  \item \textbf{Find the singular values.} Compute $A\T A$, find its eigenvalues $\lambda$, and give the
        singular values $\sigma_1,\sigma_2$ (larger first).
        \vspace{1.7cm}
  \item \textbf{One output vector.} The input singular vector for $\lambda_1$ is
        $\vv_1=\frac{1}{\sqrt5}\begin{bmatrix} 1 \\ 2 \end{bmatrix}$. Compute the output singular vector
        $\uu_1=A\vv_1/\sigma_1$.
        \vspace{1.7cm}
  \item \textbf{Explain.} Why do the eigenvalues of $A\T A$ always let us take a real
        $\sigma=\sqrt{\lambda}$? (What is special about $A\T A$, and what does that guarantee about
        its eigenvalues?)
        \writelines{2}
\end{enumerate}

\end{document}
